% problem_id: S001-lemma-smash-nilpotence
% source_id: S001 (Stacks Project)
% source_file: weil.tex
% statement_locator: weil.tex lines 2185--2198
% proof_locator: weil.tex lines 2200--2334
% env: lemma
% id_note: label 在本来源内唯一
% pairing: tight (verified)
% status: complete (statement + author proof verbatim + reason + locators)
%
\begin{lemma}[Voevodsky]
\label{lemma-smash-nilpotence}
\begin{reference}
\cite{nilpotence}
\end{reference}
Let $k$ be a field. Let $X$ be a geometrically irreducible
smooth projective scheme over $k$. Let $x, x' \in X$ be $k$-rational points.
For $n$ large enough the class of the zero cycle
$$
([x] - [x']) \times \ldots \times ([x] - [x']) \in
\CH_0(X^n)
$$
is torsion.
\end{lemma}

\begin{proof}
If we can show this after base change to the algebraic closure of $k$,
then the result follows over $k$ because the kernel of pullback
is torsion by Lemma \ref{lemma-kernel-to-closure}.
Hence we may and do assume $k$ is algebraically closed.

\medskip\noindent
Using Bertini we can choose a smooth curve $C \subset X$ passing through
$x$ and $x'$. See proof of Lemma \ref{lemma-divide-difference-points}.
Hence we may assume $X$ is a curve.

\medskip\noindent
Assume $X$ is a curve and $k$ is algebraically closed.
Write $S^n(X) = \underline{\Hilbfunctor}^n_{X/k}$ with notation as in
Picard Schemes of Curves, Sections \ref{pic-section-hilbert-scheme-points}
and \ref{pic-section-divisors}. There is a canonical morphism
$$
\pi : X^n \longrightarrow S^n(X)
$$
which sends the $k$-rational point $(x_1, \ldots, x_n)$ to the $k$-rational
point corresponding to the divisor $[x_1] + \ldots + [x_n]$ on $X$.
There is a faithful action of the symmetric group $S_n$ on $X^n$.
The morphism $\pi$ is $S_n$-invariant and the fibres of $\pi$ are
$S_n$-orbits (set theoretically). Finally, $\pi$ is finite flat of
degree $n!$, see Picard Schemes of Curves, Lemma
\ref{pic-lemma-universal-object}.

\medskip\noindent
Let $\alpha_n$ be the zero cycle on $X^n$ given by the formula in the
statement of the lemma. Let $\mathcal{L} = \mathcal{O}_X(x - x')$. Then
$c_1(\mathcal{L}) \cap [X] = [x] - [x']$. Thus
$$
\alpha_n = c_1(\mathcal{L}_1) \cap \ldots \cap c_1(\mathcal{L}_n) \cap [X^n]
$$
where $\mathcal{L}_i = \text{pr}_i^*\mathcal{L}$ and $\text{pr}_i : X^n \to X$
is the $i$th projection. By either
Divisors, Lemma \ref{divisors-lemma-finite-locally-free-has-norm} or
Divisors, Lemma \ref{divisors-lemma-norm-in-normal-case}
there is a norm for $\pi$. Set
$\mathcal{N} = \text{Norm}_\pi(\mathcal{L}_1)$, 
see Divisors, Lemma \ref{divisors-lemma-norm-invertible}. We have
$$
\pi^*\mathcal{N} =
(\mathcal{L}_1 \otimes \ldots \otimes \mathcal{L}_n)^{\otimes (n - 1)!}
$$
in $\Pic(X^n)$ by a calculation. Deails omitted; hint: this follows from
the fact that
$\text{Norm}_\pi : \pi_*\mathcal{O}_{X^n} \to \mathcal{O}_{S^n(X)}$
composed with the natural map $\pi_*\mathcal{O}_{S^n(X)} \to \mathcal{O}_{X^n}$
is equal to the product over all $\sigma \in S_n$ of the action of
$\sigma$ on $\pi_*\mathcal{O}_{X^n}$. Consider
$$
\beta_n = c_1(\mathcal{N})^n \cap [S^n(X)]
$$
in $\CH_0(S^n(X))$. Observe that
$c_1(\mathcal{L}_i) \cap c_1(\mathcal{L}_i) = 0$
because $\mathcal{L}_i$ is pulled back from a curve, see
Chow Homology, Lemma \ref{chow-lemma-vanish-above-dimension}. Thus we see that
\begin{align*}
\pi^*\beta_n
& =
((n - 1)!)^n
(\sum\nolimits_{i = 1, \ldots, n} c_1(\mathcal{L}_i))^n \cap [X^n] \\
& =
((n - 1)!)^n n^n 
c_1(\mathcal{L}_1) \cap \ldots \cap c_1(\mathcal{L}_n) \cap [X^n] \\
& =
(n!)^n \alpha_n
\end{align*}
Thus it suffices to show that $\beta_n$ is torsion.

\medskip\noindent
There is a canonical morphism
$$
f : S^n(X) \longrightarrow \underline{\Picardfunctor}^n_{X/k}
$$
See Picard Schemes of Curves, Lemma \ref{pic-lemma-picard-pieces}.
For $n \geq 2g - 1$ this morphism is a projective space bundle
(details omitted; compare with the
proof of Picard Schemes of Curves, Lemma \ref{pic-lemma-picard-pieces}).
The invertible sheaf $\mathcal{N}$ is trivial on the
fibres of $f$, see below. Thus by the projective space bundle formula
(Chow Homology, Lemma \ref{chow-lemma-chow-ring-projective-bundle})
we see that $\mathcal{N} = f^*\mathcal{M}$ for some invertible
module $\mathcal{M}$ on $\underline{\Picardfunctor}^n_{X/k}$.
Of course, then we see that
$$
c_1(\mathcal{N})^n = f^*(c_1(\mathcal{M})^n)
$$
is zero because $n > g = \dim(\underline{\Picardfunctor}^n_{X/k})$
and we can use Chow Homology, Lemma \ref{chow-lemma-vanish-above-dimension}
as before.

\medskip\noindent
We still have to show that $\mathcal{N}$ is trivial on a fibre $F$
of $f$. Since the fibres of $f$ are projective spaces and since
$\Pic(\mathbf{P}^m_k) = \mathbf{Z}$
(Divisors, Lemma \ref{divisors-lemma-Pic-projective-space-UFD}),
this can be shown by computing the degree of $\mathcal{N}$
on a line contained in the fibre. Instead we will prove it by
proving that $\mathcal{N}$ is algebraically
equivalent to zero. First we claim there is a connected finite type scheme $T$
over $k$, an invertible module $\mathcal{L}'$ on $T \times X$ and
$k$-rational points $p, q \in T$ such that
$\mathcal{M}_p \cong \mathcal{O}_X$ and $\mathcal{M}_q = \mathcal{L}$.
Namely, since $\mathcal{L} = \mathcal{O}_X(x - x')$ we can take
$T = X$, $p = x'$, $q = x$, and
$\mathcal{L}' = \mathcal{O}_{X \times X}(\Delta)
\otimes \text{pr}_2^*\mathcal{O}_X(-x')$.
Then we let $\mathcal{L}'_i$ on
$T \times X^n$ for $i = 1, \ldots, n$
be the pullback of $\mathcal{L}'$ by
$\text{id}_T \times \text{pr}_i : T \times X^n \to T \times X$.
Finally, we let
$\mathcal{N}' = \text{Norm}_{\text{id}_T \times \pi}(\mathcal{L}'_1)$
on $T \times S^n(X)$.
By construction we have $\mathcal{N}'_p = \mathcal{O}_{S^n(X)}$
and $\mathcal{N}'_q = \mathcal{N}$.
We conclude that
$$
\mathcal{N}'|_{T \times F}
$$
is an invertible module on $T \times F \cong T \times \mathbf{P}^m_k$
whose fibre over $p$ is the trivial invertible module and whose fibre
over $q$ is $\mathcal{N}|_F$. Since the euler characteristic
of the trivial bundle is $1$ and since this euler characteristic
is locally constant in families (Derived Categories of Schemes,
Lemma \ref{perfect-lemma-chi-locally-constant-geometric})
we conclude $\chi(F, \mathcal{N}^{\otimes s}|_F) = 1$
for all $s \in \mathbf{Z}$. This can happen only if
$\mathcal{N}|_F \cong \mathcal{O}_F$ (see
Cohomology of Schemes, Lemma
\ref{coherent-lemma-cohomology-projective-space-over-ring})
and the proof is complete. Some details omitted.
\end{proof}
